%==============================================================================
%  CV — Nathan Palvadeau
%  Réplique LaTeX du CV Canva "Security Monitoring Analyst — 2026"
%
%  Moteur   : XeLaTeX (obligatoire : fontspec + polices TTF locales)
%  Overleaf : Menu > Paramètres > Compilateur = XeLaTeX
%             et uploader le dossier fonts/ à la racine du projet
%  Local    : xelatex -output-directory=build cv-nathan-palvadeau.tex
%==============================================================================
\documentclass[9pt,a4paper]{extarticle}

%--- Géométrie (relevée au point près sur le PDF Canva d'origine) -------------
\usepackage[a4paper,
            left=25.5pt, right=14.3pt,
            top=25pt,    bottom=14pt,
            nohead, nofoot]{geometry}

\usepackage{fontspec}
\usepackage{xcolor}
\usepackage{enumitem}
\usepackage{changepage}
\usepackage{calc}
\usepackage[hidelinks]{hyperref}

%==============================================================================
%  COULEURS
%==============================================================================
\definecolor{paper}{HTML}{FAFAFA}   % fond de page du Canva
\definecolor{ink}{HTML}{000000}     % texte
\pagecolor{paper}
\color{ink}

%==============================================================================
%  POLICES
%  Corps  : DM Sans       (Google Fonts, OFL) — police d'origine du Canva
%  Titres : Space Grotesk (Google Fonts, OFL) — substitut libre de
%           PP Neue Machina (commerciale) utilisée par le Canva.
%           Si tu possèdes PP Neue Machina : dépose les .ttf dans fonts/
%           et remplace SpaceGrotesk ci-dessous par PPNeueMachina.
%==============================================================================
\setmainfont{DMSans}[
  Path            = ./fonts/,
  Extension       = .ttf,
  UprightFont     = *-Regular,
  BoldFont        = *-Bold,
  ItalicFont      = *-Italic,
  BoldItalicFont  = *-BoldItalic,
  Ligatures       = TeX,
]
\newfontfamily\displayfont{SpaceGrotesk}[
  Path        = ./fonts/,
  Extension   = .ttf,
  UprightFont = *-Medium,
  BoldFont    = *-Bold,
  Ligatures   = TeX,
]

%==============================================================================
%  COMPOSITION : ferré à gauche, sans césure (comme le Canva)
%==============================================================================
\raggedright
\hyphenpenalty=10000
\exhyphenpenalty=10000
\tolerance=9999
\emergencystretch=0pt
\setlength{\parindent}{0pt}
\setlength{\parskip}{0pt}
\pagestyle{empty}

%==============================================================================
%  GRILLE
%==============================================================================
\newlength{\labelcol}\setlength{\labelcol}{125.1pt}  % gouttière d'intitulés
\newlength{\colgap}  \setlength{\colgap}{222.2pt}    % pas de la grille 2 col.
\newlength{\cellw}   \setlength{\cellw}{208.2pt}     % largeur d'une cellule
\newcommand{\ruleweight}{0.75pt}

%--- Écarts verticaux (tout le rythme vertical se règle ici) ------------------
\newlength{\gNameContact}  \setlength{\gNameContact}  {7.9pt}
\newlength{\gContactRule}  \setlength{\gContactRule}  {4pt}
\newlength{\gRuleSection}  \setlength{\gRuleSection}  {1.9pt}
\newlength{\gTitleOrg}     \setlength{\gTitleOrg}     {2.0pt}
\newlength{\gOrgStack}     \setlength{\gOrgStack}     {2.5pt}
\newlength{\gStackItems}   \setlength{\gStackItems}   {6pt}
\newlength{\gJobJob}       \setlength{\gJobJob}       {10pt}
\newlength{\gItemsRule}    \setlength{\gItemsRule}    {2.9pt}
\newlength{\gSchoolDegree} \setlength{\gSchoolDegree} {3.6pt}
\newlength{\gDegreeItems}  \setlength{\gDegreeItems}  {4.0pt}
\newlength{\gGridRow}      \setlength{\gGridRow}      {2pt}
\newlength{\gCellTitleSub} \setlength{\gCellTitleSub} {3.5pt}
\newlength{\gSchoolSchool} \setlength{\gSchoolSchool}{7pt}
\newlength{\gGridRule}    \setlength{\gGridRule}    {8.8pt}

%==============================================================================
%  PRIMITIVES
%==============================================================================
% Corps de texte dimensionné : \sz{taille}{interligne}
\newcommand{\sz}[2]{\fontsize{#1}{#2}\selectfont}

% Intitulé de section, posé dans la gouttière de gauche.
% La boîte ne doit contribuer ni hauteur ni profondeur, sinon elle écarte la
% ligne suivante : on l'aplatit avec \raisebox (\smash n'accepte pas de \par).
\newcommand{\flatbox}[1]{\raisebox{0pt}[0pt][0pt]{#1}}
\newcommand{\sidelabel}[1]{%
  \leavevmode
  \llap{\flatbox{\parbox[t]{\labelcol}{\raggedright\bfseries\sz{10}{13.5}#1}}}%
  \ignorespaces
}

% Filet horizontal, pleine largeur de la colonne contenu
\newcommand{\cvrule}{%
  \par\noindent\textcolor{ink}{\rule{\linewidth}{\ruleweight}}\par%
}

% Puce de largeur nulle : le texte reste aligné sur \leftmargin
\newcommand{\bul}{\makebox[0pt][l]{\sz{10}{10}\textbullet}}

% Flèche de la ligne « Tech Stack »
\newcommand{\arrow}{→}

\newcommand{\techstack}[1]{%
  {\sz{8}{11}Tech Stack\ \arrow\ \textbf{#1}\par}%
}

% En-tête d'expérience : poste / organisation / stack technique
\newcommand{\job}[3]{%
  {\sz{10}{13.5}\bfseries\leavevmode #1\par}%
  \vspace{\gTitleOrg}%
  {\sz{8.9}{12}#2\par}%
  \vspace{\gOrgStack}%
  \techstack{#3}%
}

% Puces d'expérience : texte à +9.5pt, puce à -2.1pt de la colonne
\newenvironment{cvitems}{%
  \vspace{\gStackItems}%
  \begin{itemize}[leftmargin=9.5pt, labelsep=11.6pt, labelwidth=0pt,
                  align=left, itemsep=0pt, topsep=0pt, parsep=0pt,
                  partopsep=0pt, label=\bul]%
  \sz{9}{12.5}%
}{\end{itemize}}

% Puces de formation : texte à +16pt, puce à +4.5pt
\newenvironment{cvsubitems}{%
  \vspace{\gDegreeItems}%
  \begin{itemize}[leftmargin=16pt, labelsep=11.5pt, labelwidth=0pt,
                  align=left, itemsep=0pt, topsep=0pt, parsep=0pt,
                  partopsep=0pt, label=\bul]%
  \sz{9}{12.5}%
}{\end{itemize}}

% Cellule de la grille « Projets » / « Langues »
\newcommand{\gridcell}[2]{%
  \parbox[t]{\cellw}{%
    \sz{9.05}{12}\raggedright
    \textbf{#1}\par\vspace{\gCellTitleSub}#2\par}%
}
% Rangée de deux cellules
\newcommand{\gridrow}[4]{%
  \gridcell{#1}{#2}\hspace*{14pt}\gridcell{#3}{#4}%
}

\begin{document}

%==============================================================================
%  EN-TÊTE
%==============================================================================
\vspace*{-4pt}

\begin{adjustwidth}{0pt}{11.7pt}

% Nom (gauche) + intitulé de poste (droite), lignes de base alignées
\noindent
{\displayfont\bfseries\sz{23}{23}Nathan PALVADEAU}%
\hfill
{\displayfont\sz{14}{14}\addfontfeature{LetterSpace=4.0}Security Monitoring Analyst}%
\par

\vspace{\gNameContact}

% Coordonnées (gauche) + disponibilité (droite)
\noindent
\parbox[t]{295pt}{%
  \sz{9}{12}\raggedright
  \textbf{nathan.palvadeau.pro@proton.me} \textbullet{}
  \href{https://linkedin.com/in/nathan-plvd}{\underline{linkedin.com/in/nathan-plvd}}\par
  \textbf{+33 7 82 81 73 64} \textbullet{} Versailles (78) \textbullet{} Permis B\par
}%
\hfill
\parbox[t]{235pt}{%
  \sz{8.75}{11.7}\raggedleft
  \textbf{Recherche alternance sur 2 ans}\par
  4 semaines entreprise / 2 semaines école\par
}

\end{adjustwidth}

%==============================================================================
%  CORPS — gouttière d'intitulés + colonne de contenu
%==============================================================================
\begin{adjustwidth}{\labelcol}{0pt}

\vspace{\gContactRule}
\cvrule
\vspace{\gRuleSection}

%------------------------------------------------------------------------------
%  PROFIL
%------------------------------------------------------------------------------
\sidelabel{PROFIL}%
{\sz{9}{13.0}\leavevmode Analyste SOC N2 en alternance sur un périmètre européen
anglophone de +120k endpoints. Parcours infrastructure (bachelor systèmes et
réseaux) puis cybersécurité. Investigation, ingénierie de détection et
automatisation SOC, nourries par une pratique offensive continue (CRTP,
plateformes CTF).\par}

\vspace{\gItemsRule}
\cvrule
\vspace{\gRuleSection}

%------------------------------------------------------------------------------
%  EXPÉRIENCES PROFESSIONNELLES
%------------------------------------------------------------------------------
\sidelabel{EXPÉRIENCES PROFESSIONNELLES}%
\job{Analyste SOC N2 \textbullet{} Scope Européen (anglophone) \textbullet{} +120k endpoints}
    {Cyna (Paris, FR) | 7 analystes | \textit{Depuis août 2025}}
    {Elastic, Microsoft Defender, SentinelOne, TheHive \& Cortex, n8n, Zendesk, Python, Claude Code}

\begin{cvitems}
  \item \textbf{Analyse et qualification} des incidents de sécurité : phishing O365,
        compromission de postes, mouvements latéraux, Active Directory, attaques par
        force brute, infections malware et enrichissement contextuel pour appuyer les
        investigations.
  \item Réalisation d'\textbf{investigations approfondies}, mise en œuvre des actions
        de confinement et coordination technique avec les clients et les éditeurs de
        solutions de sécurité.
  \item Amélioration des \textbf{capacités de détection} : création de règles de
        détection avancées, automatisation de playbooks et activités de threat hunting.
  \item Contribution au développement des capacités SOC : formation, documentation,
        standardisation de procédures et pilotage de projet interne.
  \item Collaboration avec l'équipe \textbf{R\&D} : prototypage, tests et déploiement
        d'améliorations liées à l'\textbf{automatisation SOC} et à la détection
        (intégration IA, \textbf{CTI}, règles firewall, corrélation de logs).
\end{cvitems}

\vspace{\gJobJob}

\job{Cyber Forensic \textbullet{} Entreprise internationale \textbullet{} +10k utilisateurs}
    {ADP (Nanterre, FR) | 4 analystes | \textit{Août 2024}}
    {CrowdStrike, Microsoft Defender, E-Discovery, Volatility, Autopsy, Python, PowerShell, OpenAI}

\begin{cvitems}
  \item \textbf{Forensic} et \textbf{réponse à incident} : analyses approfondies de
        malwares et investigation d'artefacts afin de reconstituer des scénarios
        d'attaque complexes.
  \item Opérations DFIR et workflows d'investigation : utilisation d'outils spécialisés
        pour la collecte, la préservation et l'\textbf{analyse d'artefacts}, avec
        automatisation des étapes récurrentes.
  \item Environnement de recherche dédié : conception d'un laboratoire d'investigation
        \textbf{OSINT} (détails omis pour des raisons de confidentialité) et rédaction
        de la documentation associée.
  \item Pilotage d'un \textbf{projet à l'international} : déploiement de
        l'environnement dans le cadre d'une investigation réelle et présentation des
        conclusions stratégiques à une équipe technique basée aux États-Unis.
\end{cvitems}

\vspace{\gItemsRule}
\cvrule
\vspace{\gRuleSection}

%------------------------------------------------------------------------------
%  FORMATION
%------------------------------------------------------------------------------
\sidelabel{FORMATION}%
{\sz{10}{13.5}\bfseries\leavevmode Ecole2600\par}
\vspace{\gSchoolDegree}
{\sz{9.76}{13}Mastère en cybersécurité (programme en alternance)\par}

\begin{cvsubitems}
  \item Sécurité applicative \& offensive : sécurité Web, tests d'intrusion et
        évaluation de vulnérabilités.
  \item Reverse engineering, exploitation bas niveau (pwn) \& analyse malware
  \item Sécurité infrastructure \& cloud : réseaux, systèmes, sécurité cloud et
        pratiques DevSecOps.
  \item Threat Intelligence \& recherche appliquée : Cyber Threat Intelligence (CTI),
        gestion de projet et recherches spécialisées.
  \item Gouvernance, risque et compliance
  \item Programmation : C, Python \& Rust
\end{cvsubitems}

\vspace{\gSchoolSchool}

{\sz{10}{13.5}\bfseries\leavevmode Sup de Vinci\par}
\vspace{\gSchoolDegree}
{\sz{9.76}{13}Bachelor Administration Systèmes \& Réseaux (2 ans)\par}

\vspace{\gItemsRule}
\cvrule
\vspace{\gRuleSection}

%------------------------------------------------------------------------------
%  PROJETS
%------------------------------------------------------------------------------
\sidelabel{PROJETS}%
\gridrow{Certified Red Team Professional (CRTP)}{Altered Security - En cours (2026)}
        {CTF Platforms}{Autodidacte | Depuis 2023\\HackTheBox, Root-Me, TryHackMe}

\vspace{\gGridRow}

\noindent
\gridrow{Blockchain \& Automatisations}{Projet de R\&D sur des thématiques de FinInt, tracking de transactions décentralisées}
        {SOC \& Home Lab Infrastructure}{Autodidacte | Depuis 2023}

\vspace{\gGridRule}
\cvrule
\vspace{\gRuleSection}

%------------------------------------------------------------------------------
%  LANGUES
%------------------------------------------------------------------------------
\sidelabel{LANGUES}%
\gridrow{Français}{Natif}{Anglais}{Capacité professionnelle complète (B2)}

\end{adjustwidth}
\end{document}
